% !TeX program = xelatex
% !BIB program = biber
\documentclass[light]{lutbeamer}
\usepackage{pgfpages}
\usepackage{ragged2e}
\usepackage{booktabs}
\usepackage{amssymb}
\usepackage{fontawesome5}
\usepackage{tikz}
\usepackage{pgfplots}
\pgfplotsset{compat=1.18}
\addbibresource{references.bib}

% ===================================================================
% METADATA
% Replace placeholders after copying this asset into the target week file.
% ===================================================================
\setdepartment{OSTIM Teknik Üniversitesi}
\institute{}
\author[Mühendislik Fakültesi]{\\
Prof.~Dr.~Yalçın ATA \\
Prof.~Dr.~Arif DEMİR \\
Dr.~Öğr.~Üyesi~Muhammed ELMNEFİ \\
Dr.~Öğr.~Üyesi~Haydar KILIÇ \\
Dr.~Öğr.~Üyesi~Emel GÜVEN \\
Dr.~Öğr.~Üyesi~Ufuk ASIL \\
Öğr.~Gör.~Sema ÇİFTÇİ
}
\title{<COURSE_TITLE>}
\subtitle{Hafta <WEEK_NO>: <TOPIC>}
\date{2025--2026 Bahar Dönemi}

% ===================================================================
% AUTOMATIC TABLE OF CONTENTS
% Tune the section split after the real section count is known.
% ===================================================================
\setcounter{tocdepth}{2}

\AtBeginSection[]{
  \begin{frame}{Sunum Planı}
    \scriptsize
    \begin{columns}[T]
      \begin{column}{0.48\textwidth}
        \tableofcontents[sections={1-4}, currentsection]
      \end{column}
      \begin{column}{0.48\textwidth}
        \tableofcontents[sections={5-10}, currentsection]
      \end{column}
    \end{columns}
  \end{frame}
}

\begin{document}

\setbeamertemplate{navigation symbols}{}
\setbeamertemplate{footline}{
  \begin{beamercolorbox}[wd=\paperwidth, ht=10pt, dp=1pt]{footline}
    \usebeamerfont{footline}
    \begin{tikzpicture}[remember picture,overlay]
      \node[anchor=south west, xshift=10mm, yshift=.4mm]
      at (current page.south west)
      {\insertshortauthor,~\department};
    \end{tikzpicture}
    \begin{tikzpicture}[remember picture,overlay]
      \node[anchor=south east, xshift=-5mm, yshift=0.5mm]
      at (current page.south east)
      {\insertframenumber \,/\, \inserttotalframenumber};
    \end{tikzpicture}
  \end{beamercolorbox}
}

% ===================================================================
% COVER
% ===================================================================
\begin{frame}<article:0>[plain,noframenumbering]
  \begin{tikzpicture}[remember picture,overlay]
    \node[anchor=center, at=(current page.center), inner sep=0pt] {
      \includegraphics[width=\paperwidth, height=\paperheight]{logos/background_figure.jpg}
    };
  \end{tikzpicture}
\end{frame}

{
  \setbeamertemplate{footline}{
    \begin{beamercolorbox}[wd=\paperwidth, ht=10pt, dp=1pt]{footline}
      \usebeamerfont{footline}
      \begin{tikzpicture}[remember picture,overlay]
        \node[anchor=south west, xshift=10mm, yshift=.4mm, align=left]
        at (current page.south west)
        {\tiny Bu şablon OTU Yazılım Mühendisliği tarafından hazırlanmış olup \href{https://github.com/ufukasia/Ostim-Teknik--niversitesi-Yaz-l-m-M-hendisli-i---Lutbeamer-egitim-i-eri-i-haz-rlama-}{buradan} güncel haline ulaşabilirsiniz.};
      \end{tikzpicture}
      \begin{tikzpicture}[remember picture,overlay]
        \node[anchor=south east, xshift=-5mm, yshift=0.5mm]
        at (current page.south east)
        {\insertframenumber \,/\, \inserttotalframenumber};
      \end{tikzpicture}
    \end{beamercolorbox}
  }
  \maketitle{}
}

\begin{frame}{İçindekiler}
  \scriptsize
  \begin{columns}[T]
    \begin{column}{0.48\textwidth}
      \tableofcontents[sections={1-4}]
    \end{column}
    \begin{column}{0.48\textwidth}
      \tableofcontents[sections={5-10}]
    \end{column}
  \end{columns}
\end{frame}

% ===================================================================
% CONTENT START
% Türkçe içerikte ASCII transliterasyon kullanma; gerçek Türkçe karakterleri koru.
% Fill section / subsection / frame titles first, then content.
% ===================================================================

% \section{Köprü veya Motivasyon}
% \subsection{Hatirlatma}
% \begin{frame}{Hafta N'den Ne Getiriyoruz?}
% \end{frame}

% \section{Ana Konu 1}
% \subsection{Temel Kavram}
% \begin{frame}{Kavrama Giriş}
% \end{frame}
% \begin{frame}{Sayisal Örnek}
% \end{frame}
% \subsection{Pekiştirme Sorusu}
% \begin{frame}[t]{Soru 1: ...}
% \end{frame}
% \begin{frame}[t]{Soru 1: Çözüm İçin Durak}
% \end{frame}
% \begin{frame}[t]{Soru 1: Çözüm}
% \end{frame}

% ===================================================================
% CONTENT END
% ===================================================================

\section{Kaynaklar}
\begin{frame}[allowframebreaks]{Kaynaklar}
  \printbibliography{}
\end{frame}

\appendix

\section{Yapay Zeka Asistanı}
\begin{frame}{Yapay Zeka Asistanınız: NotebookLM}
  \begin{columns}[T]
    \begin{column}{0.48\textwidth}
      \justifying{}
      \textbf{Ders Materyalleri ile Etkilesime Gecin!}

      Bu dersin sunumlari, notlari ve yardimci icerikleri Google'in gelismis yapay zeka Asistanı \textbf{NotebookLM} sisteminde duzenlenebilir.

      \vspace{1em}
      \begin{itemize}
        \item \textbf{Aninda Cevaplar:} Akliniza takilanlari sorun.
        \item \textbf{Özetler:} Karmasık konulari hizla toparlayin.
        \item \textbf{Audio Overview:} Dersi podcast gibi dinleyin.
      \end{itemize}
    \end{column}
    \begin{column}{0.48\textwidth}
      \centering
      \includegraphics[width=0.82\linewidth]{figures/notebooklm.jpg}
    \end{column}
  \end{columns}
\end{frame}

\begin{frame}[plain,noframenumbering]
  \begin{tikzpicture}[remember picture,overlay]
    \node[anchor=center, at=(current page.center), inner sep=0pt] {
      \includegraphics[width=\paperwidth, height=\paperheight]{logos/background_figure.jpg}
    };
  \end{tikzpicture}
\end{frame}

\end{document}

